从医疗容量选择阈值与检验干预能力是两个步骤。若按上述平均需求充分条件选取 $\eta$，还须满足下文的触发条件；有额外能力限制时，再分别检查仅隔离或联合控制的实施判据。满足这些条件后，同一医疗需求上界适用于所比较的全部允许分配，但不意味着它们具有相同的实际需求轨迹。后文只在共同启动状态、阈值及退出目标下比较控制，不把这一充分条件作为新的时变状态约束加入成本最小化。

\subsection{常规传播轨迹与触发条件}
